\documentclass[11pt,a4paper]{article}
\usepackage[italian]{babel}
\usepackage[margin=2.5cm]{geometry}
\usepackage{parskip}
\usepackage{amsmath, amssymb, amsthm}
\usepackage[most]{tcolorbox}
\usepackage{tikz}
\usepackage{hyperref}

% --- Riquadri con bordo nero, senza numero ---
\tcbset{nota/.style={colback=white, colframe=black, boxrule=0.5pt,
  sharp corners}}
\NewTColorBox{definizione}{}{nota,
  before upper={\textbf{Definizione.}\ }}
\NewTColorBox{teorema}{o}{nota,
  before upper={\textbf{Teorema\IfValueT{#1}{ (#1)}.}\ }}
\NewTColorBox{proposizione}{o}{nota,
  before upper={\textbf{Proposizione\IfValueT{#1}{ (#1)}.}\ }}

% --- Ambienti semplici (senza riquadro) ---
\theoremstyle{definition}
\newtheorem*{esempio}{Esempio}
\newtheorem*{osservazione}{Osservazione}

\title{Funzioni: la relazione associata}
\author{Marco Cerbella}
\date{Lezione 4 -- 29 settembre 2026}

\begin{document}
\maketitle

Si riprendono le funzioni introdotte nella lezione 3 (definizione, dominio, codominio, grafico, funzioni costanti).

\section{La relazione associata a una funzione}

\begin{definizione}
Sia $f : A \to B$ una funzione. La \emph{relazione associata a $f$}, indicata
con $R_f$, è la relazione su $A$ definita da
\[
a_1 \mathrel{R_f} a_2 \iff f(a_1) = f(a_2).
\]
\end{definizione}

Due elementi di $A$ sono quindi in relazione se $f$ associa loro la stessa
immagine.

\begin{esempio}
Siano $A = \{1, 2, 3, 4\}$ e $B = \{x, y, z\}$ e sia $f : A \to B$ data da
\[
f(1) = f(2) = f(3) = x, \qquad f(4) = y.
\]
Allora $1 \mathrel{R_f} 2$, $3 \mathrel{R_f} 3$, $4 \mathrel{R_f} 4$, \dots{}
Le classi di equivalenza sono
\[
[1]_{R_f} = \{1, 2, 3\} = [2]_{R_f} = [3]_{R_f}, \qquad
[4]_{R_f} = \{4\},
\]
e l'insieme quotiente è $A/_{R_f} = \{[1]_{R_f}, [4]_{R_f}\}$.
\end{esempio}

\begin{teorema}
Sia $f : A \to B$ una funzione. La relazione $R_f$ è una relazione di
equivalenza su $A$.
\end{teorema}

\begin{proof}
Lasciata per esercizio a lezione; la verifica è
immediata, perché ``avere la stessa immagine'' eredita le proprietà
dell'uguaglianza in $B$.
\begin{enumerate}
    \item \emph{Riflessiva}: per ogni $a \in A$ vale $f(a) = f(a)$, quindi
          $a \mathrel{R_f} a$.
    \item \emph{Simmetrica}: se $a_1 \mathrel{R_f} a_2$, cioè
          $f(a_1) = f(a_2)$, allora $f(a_2) = f(a_1)$, cioè
          $a_2 \mathrel{R_f} a_1$.
    \item \emph{Transitiva}: se $a_1 \mathrel{R_f} a_2$ e
          $a_2 \mathrel{R_f} a_3$, allora $f(a_1) = f(a_2)$ e
          $f(a_2) = f(a_3)$, quindi $f(a_1) = f(a_3)$, cioè
          $a_1 \mathrel{R_f} a_3$. \qedhere
\end{enumerate}
\end{proof}

\end{document}
